Understanding why ``agentic AI for software testing'' has become a
distinct and fast-moving research area requires situating it against
the longer trajectory of test-automation research, clarifying the
specific technical properties of large language models (LLMs) and
agentic architectures that differentiate this generation of tools from
earlier automation, and fixing a shared vocabulary for terms --
autonomy level, self-healing, the test-oracle problem, flakiness --
that recur throughout the taxonomy and synthesis in later sections.
This section does not attempt a comparative analysis of prior surveys
(that is the task of Section~\ref{sec:related-work}); instead, it
supplies the conceptual scaffolding a reader needs before engaging with
the review's results.

\subsection{A Short History of Software Test Automation}
Software testing automation has evolved through several overlapping
eras, each lowering the manual effort required to exercise a system
while raising new questions about what ``coverage'' or ``correctness''
actually guarantees. The earliest automation targeted the mechanics of
test \emph{execution} rather than test \emph{design}: capture/replay
tools recorded a tester's interactions with a GUI or API and replayed
them as scripted regression suites. This approach is cheap to adopt but
notoriously brittle -- a single UI change invalidates every recorded
locator, producing the maintenance burden that decades later motivates
self-healing locator research~\cite{h2nd,rohamo2024robotframework}.
A second era shifted the unit of automation from whole user journeys to
isolated code units, popularized by xUnit-style frameworks that made
writing and running assertions over individual functions or classes a
routine developer activity rather than a separate QA discipline; this
coverage-driven, developer-authored style of unit testing remains the
substrate that most LLM-based test-generation work still targets and
evaluates against~\cite{slr0072025,slr0082025}.

A third era pursued automated test \emph{generation} rather than mere
execution, using the structure of the program itself as a generation
signal. Search-based software testing framed test-input selection as
an optimization problem over a fitness landscape (typically branch or
line coverage), while symbolic and concolic execution reasoned about
path constraints to derive inputs that reach specific program
locations; fuzzing complemented both by mutating or generating inputs
to find crashes and robustness violations without requiring a precise
specification of correctness. These coverage-driven and property-driven
generators were effective at producing inputs but persistently weak at
producing \emph{meaningful} test oracles and at escaping coverage
plateaus once a program's branching structure became deep or
semantically opaque -- a limitation that later motivated hybridizing
search-based generation with LLMs, as in CodaMosa's strategy of
invoking a pre-trained language model to escape stalled coverage
exploration~\cite{slr0122023}. Classical fuzzers and symbolic engines
also had no access to natural-language context (requirements,
identifiers, comments, commit messages) and so could not connect a
generated input to a domain-level notion of intended
behavior~\cite{xia2024,deng2023,deng2023b}.

The most recent era layers machine learning, and subsequently large
language models and agentic orchestration, on top of this stack.
Early ML-for-testing work used learned models for narrower sub-problems
-- prioritizing test cases, predicting flaky outcomes, or ranking
likely-buggy locations -- without altering the fundamentally
human-authored nature of test design. The arrival of LLMs trained on
large code and natural-language corpora changed the unit of automation
again: a model can now read a function, its surrounding repository
context, and an informal description of intent, and directly emit a
syntactically valid, semantically plausible unit test, mutation, fuzz
target, or GUI action sequence~\cite{slr0012024,slr0022024,slr0032023}.
Agentic AI represents a further step beyond single-shot generation: a
testing agent couples an LLM with tool use (running a compiler,
executing a test, reading a stack trace, querying a repository) inside
an iterative plan-act-observe loop, allowing it to generate a test,
execute it, diagnose a failure, and revise its own output without
further human intervention~\cite{slr0232024,slr0252024,slr0242024}.
Multi-agent variants extend this further still, assigning specialized
roles (generator, critic, oracle-repairer) to separate LLM instances
that debate or vote on a shared artifact before it is
accepted~\cite{slr0192025,xu2025candor,slr0222026}. This progression --
from replaying what a human already did, to generating code that
satisfies a structural criterion, to agents that plan, act, and revise
toward a goal -- is the backdrop against which the remainder of this
review classifies next-generation testing techniques.

\subsection{Why LLMs and Agentic AI Changed the Picture}
Four technical properties of LLM-based and agentic systems explain why
this generation of tools is treated as a qualitative departure from
earlier automation rather than an incremental improvement to it.

\textbf{In-context code understanding.} Unlike symbolic or
coverage-driven generators, which reason purely over a program's
syntax or control-flow graph, LLMs are trained on vast corpora of code
paired with natural-language artifacts (comments, documentation,
commit messages, issue discussions), giving them an implicit model of
what a function is \emph{supposed} to do, not merely what paths it
contains. This lets a model propose plausible test oracles -- expected
return values, exception conditions, invariants -- directly from a
function's signature, docstring, and surrounding usage, something
earlier generators could only approximate through heuristics such as
implicit specifications mined from existing assertions or crash
oracles~\cite{dinella2022,hossain2023,o32024}.

\textbf{Natural-language-to-test generation.} Because LLMs operate
over natural language and code in a shared representation, a
requirement, bug report, or user story can itself become the
generation prompt, collapsing the traditional pipeline in which a
human first translates a requirement into a formal test specification.
Tools such as RESTGPT for API testing, NaviQAte and WebTestPilot for
web navigation, and DroidAgent and GPTDroid for mobile GUI exploration
all exploit this property to drive functional test generation directly
from natural-language intent or app-state descriptions rather than
from pre-written scripts~\cite{kim2024,shahbandeh2024,g62026,yoon2023,liu2024}.

\textbf{Tool use and multi-step, multi-agent planning loops.} Agentic
frameworks equip an LLM with callable tools -- a shell, a debugger, a
test runner, a code-search index -- and an explicit control loop that
interleaves reasoning with action, observation, and re-planning. This
allows an agent to validate its own output empirically (does the
generated test compile? does it pass? does it reproduce the reported
bug?) rather than emitting a single unverified artifact, and to
decompose a task such as ``reproduce and fix this failing CI run''
across multiple reasoning steps and, increasingly, multiple
specialized agents~\cite{slr0232024,slr0282024,slr0602026,stennett2025}.
This loop is the architectural feature that most clearly separates
agentic testing tools from earlier one-shot neural generators, and it
is also the source of the review's later distinctions between
single-agent, tool-augmented, and multi-agent architectures (RQ2).

\textbf{Self-healing and self-adaptive test maintenance.} Because an
LLM can inspect a changed UI, API schema, or DOM tree and propose an
updated locator, assertion, or oracle in place of a now-broken one,
test \emph{maintenance} -- traditionally the single largest recurring
cost of scripted and record/replay automation -- becomes an operation
the system can perform on itself. Self-healing locator frameworks for
tools such as Robot Framework and Selenium, and runtime-testing work on
self-adaptive systems more broadly, exemplify this
shift~\cite{h2nd,rohamo2024robotframework,ijam2025healenium,joseph2026zerocost,weyns2019,wohlrab2021}.
Taken together, these four properties explain both the enthusiasm
surrounding agentic testing and the specific new risks it introduces,
since a system that can rewrite its own oracles and locators can also
silently rewrite away a genuine regression.

\subsection{Key Terminology and Concepts}
The results and discussion sections rely on a small set of terms that
are used inconsistently across the primary literature; this review
adopts the following working definitions.

\textbf{Agentic AI / autonomous testing agent.} We use \emph{agentic
AI} to denote an LLM-based system that perceives a software artifact
or execution environment, forms a plan, invokes external tools to act
on that environment, observes the result, and iterates -- as opposed to
a single-pass model invocation that maps an input prompt directly to
an output artifact~\cite{slr0042024,slr0052024,slr0062024}. An
\emph{autonomous testing agent} is an agentic system whose task is
specifically test design, generation, execution, triage, or
maintenance, potentially spanning several of these life-cycle stages
without human intervention between them~\cite{slr0882020,slr0662026}.

\textbf{Self-healing test automation.} This refers to mechanisms --
typically, but not exclusively, LLM-based -- that automatically detect
a broken element of a test (a UI locator, an API contract assumption,
an assertion that no longer matches intended behavior) and repair it at
or near runtime, without requiring a human to re-author the affected
test~\cite{h2nd,rohamo2024robotframework,ijam2025healenium}. Self-healing
is a special case of the broader notion of a self-adaptive system,
which continuously monitors and adjusts its own structure or behavior
in response to a changing environment~\cite{weyns2019,hnsel2018,sa42021}.

\textbf{The test-oracle problem.} The test-oracle problem is the
long-standing difficulty of deciding, for a given test input, what the
\emph{correct} output or behavior should be, absent a complete and
trustworthy specification. Neural and LLM-based oracle generation
(e.g., TOGA, TOGLL, AugmenTest) attempts to approximate this judgment
by learning from existing assertions, exception patterns, or developer
intent expressed in code and
comments~\cite{dinella2022,o32024,o42025,bodicoat2026oracle,molina2024oracleautomation}.
Using an LLM to both generate a test and judge whether its own
generated behavior is correct raises a circularity concern -- an agent
testing code it (or a closely related model) produced -- that is
compounded by documented self-preference biases in LLM-as-judge
settings and by evidence that models struggle to self-correct their own
reasoning or code without external
feedback~\cite{wataoka2024selfpreference,panickssery2024evaluators,huang2023selfcorrect,olausson2023selfrepair}.

\textbf{Flaky tests.} A flaky test is one that can pass or fail across
repeated executions of the same code under test, without any change to
that code, typically due to non-determinism such as timing, concurrency,
test-order dependence, or external-resource instability. Flaky tests
undermine the trustworthiness of continuous-integration signals and
have motivated a substantial body of machine-learning and, more
recently, LLM-based classification and repair
work~\cite{fatima2025,alshammari2024flakyclassifiers,parry2023cannier,leinen2026detectedundetected,parry2025systemic,lam2020lifecycle,meta2020pfs}.
Flakiness is directly relevant to self-healing test automation because
a system that automatically repairs a ``failing'' test must first
distinguish a genuine regression from an artifact of flakiness, or risk
silently suppressing real defects.

\textbf{Autonomy-level framing (L0--L5).} To compare how much human
oversight a given technique still requires, this review adopts the
autonomy-level scale used later in the roadmap and discussion
sections, following the general levels-of-autonomy framing for AI
agents of \cite{feng2025autonomylevels}: autonomy ranges from no
agentic behavior (a human performs the task with no AI assistance)
through assistive and co-pilot levels (the system suggests artifacts a
human reviews and approves) to conditional and high autonomy (the
system executes multi-step plans with narrow human checkpoints) and,
at the upper end, full autonomy, in which a testing agent designs,
executes, triages, and maintains tests across a life cycle with no
routine human intervention. This scale, together with the ISO/IEC
guidance on testing AI-based systems and on testing AI more
generally~\cite{iso2020tr29119,iso2025ts42119}, gives the taxonomy in
Section~\ref{sec:results} a consistent axis along which to classify
otherwise heterogeneous techniques (RQ2).

\subsection{Motivation: Why This Topic Needs a Rigorous Synthesis Now}
Four converging pressures motivate a systematic, rather than
narrative, review of this literature at this particular moment.

\textbf{Fragmentation of claims and terminology.} The corpus assembled
for this review spans 149 candidate studies published between 2018 and
2026, distributed across at least eighteen partially overlapping
taxonomy categories -- unit-test generation, GUI and mobile testing
agents, API testing, fuzzing, mutation testing, oracle generation,
flaky-test classification, self-adaptive runtime testing, CI-failure
triage, security of coding and testing agents, and more. Authors
routinely use different labels (``intelligent testing,'' ``autonomous
testing,'' ``AI-augmented testing,'' ``agentic software engineering'')
for closely related ideas, and several existing surveys already
attempt partial consolidations of this
space~\cite{slr0012024,slr0022024,slr0032023,slr0092024,slr0102024},
underscoring that the field has grown too large and too fast for any
single narrative account to track consistently; a reproducible,
protocol-driven synthesis is needed to make comparisons across these
overlapping categories meaningful (RQ1).

\textbf{Inconsistent benchmarks and limited generalizability.}
Empirical claims in this area are frequently anchored to small,
language-specific, or project-specific benchmarks, and several
targeted replications show that headline results are sensitive to
exactly how a benchmark is constructed and evaluated. Independent
re-evaluations of flaky-test classifiers, for instance, have found that
reported accuracy drops substantially once unrealistic class balance or
weak validation protocols are
corrected~\cite{alshammari2024flakyclassifiers,parry2023cannier}, and
studies of neural test-oracle generation report similarly large gaps
between original and large-scale replicated
evaluations~\cite{hossain2023}. Without a systematic review that
extracts and compares evaluation methodology across studies -- not
merely headline numbers -- the field risks treating isolated,
possibly fragile results as generalizable evidence (RQ3).

\textbf{Data contamination and benchmark leakage.} Because LLMs are
trained on enormous, largely undisclosed corpora that plausibly include
public GitHub repositories, there is a concrete risk that a model
``generating'' a correct test for a benchmark problem is in fact
recalling memorized training data rather than reasoning about the code
under test. Direct studies of this risk -- measuring whether LLMs
memorize widely used bug benchmarks and quantifying inter-dataset code
duplication and leakage across popular training and evaluation
corpora -- show the problem is non-hypothetical, and have motivated
the construction of contamination-resistant benchmarks built from
recently created bug
datasets~\cite{slr0302025,slr0322024,slr0312023,slr0292025}. A
synthesis that explicitly tracks which primary studies address, ignore,
or are vulnerable to contamination is necessary before the field's
quantitative claims about LLM testing effectiveness can be taken at
face value.

\textbf{The hype-versus-evidence gap and the quality-engineering safety
angle.} Enthusiastic framings of agentic AI as a wholesale replacement
for human test engineering coexist with sober industrial reports that
document narrower, assistive deployments and persistent practitioner
skepticism~\cite{slr0372025,slr0382024,slr0392025,slr0832025}, and with
emerging survey and position work explicitly interrogating how far
agentic autonomy in software engineering has actually
progressed~\cite{slr0662026,slr0692026,slr0722026,slr0842025,slr0852025}.
Closing this gap matters beyond academic tidiness because autonomous
testing agents introduce safety concerns that earlier, purely
assistive tools did not: when an agent both writes production code and
generates the tests meant to validate it, the two artifacts can share
the same blind spots, producing a circularity in which a test suite
systematically fails to catch the very class of defect the generating
model is prone to introduce~\cite{xu2025candor,slr0192025}. Testing
agents are also newly exposed, tool-using software themselves, and are
therefore subject to indirect prompt injection, unsafe tool
invocation, and sandbox-escape attacks that have been demonstrated
against LLM coding and tool-augmented agents more
broadly~\cite{zhan2024injecagent,debenedetti2024agentdojo,guo2024redcode,she2025queryipi,sakib2026trustverify,kozak2025developeraid,shi2025progent,wu2024isolategpt}.
Finally, as autonomy increases toward the higher end of the L0--L5
scale, questions of accountability, overreliance, and governance --
who is responsible when a self-healing agent masks a regression, and
what standards should govern a Level-5 testing agent -- move from
speculative to
urgent~\cite{treude2026accountable,cheng2026threepillar,passi2025overreliance,ibrahim2025overreliance,kiran2026humanai}.
Together, fragmentation, benchmark fragility, contamination risk, and
this emerging safety angle constitute the gap this systematic review
is designed to address, following an SLR protocol in the tradition of
\cite{kitchenham2007} and reported per the PRISMA
guidelines~\cite{page2021prisma}.
